\section*{Chapter \ref{ch:b3:microbiomes} --- Microbiomes and Symbioses}

\begin{solution}{exo:b3:microbiomes:1}
\omterm{def:b3:microbiomes:symbiosis}{Mutualism}: both gain --- legume and rhizobium, coral and alga.
\omterm{def:b3:microbiomes:symbiosis}{Commensalism}: one gains, the other unaffected --- most gut bacteria
feeding on what the host cannot digest. Parasitism: one gains at the
other's cost --- \emph{Wolbachia} killing male embryos. The same pair
can shift: a gut commensal that crosses the wall becomes a pathogen, a
rhizobium that stops fixing becomes a parasite of the nodule, and an
alga under heat stress harms the coral that housed it.
\end{solution}

\begin{solution}{exo:b3:microbiomes:2}
$H = -(0.5\ln 0.5 + 0.3\ln 0.3 + 0.1\ln 0.1 + 2\times 0.05\ln 0.05) =
0.347 + 0.361 + 0.230 + 0.300 = 1.24$; $e^{H} = 3.4$ effective
species. $D = 0.25 + 0.09 + 0.01 + 0.0025 + 0.0025 = 0.355$; $1/D =
2.8$.
\end{solution}

\begin{solution}{exo:b3:microbiomes:3}
The 16S survey reports who is there, to about the genus, and their
relative read counts; it says nothing about the genes they carry,
misses fungi and viruses (other \omterm{def:b3:genetic-engineering:pcr}{primers}), and is biased by copy number
and \omterm{def:b3:genetic-engineering:pcr}{primers}. Shotgun metagenomics reports genes, functions and, with
enough depth, whole genomes and strains; it costs more per sample, is
swamped by host DNA in tissue samples, and still cannot say which genes
are expressed or which cells are alive.
\end{solution}

\begin{solution}{exo:b3:microbiomes:4}
Fermenting fibre to \omterm{def:b3:microbiomes:gut}{short-chain fatty acids} that feed the colon and
the host; synthesising vitamin~K and B vitamins; transforming bile
acids and drugs; \omterm{def:b3:microbiomes:gut}{colonisation resistance} against pathogens; educating
the immune system. Disruption: \emph{Clostridioides difficile} colitis
after \omterm{def:b3:bacteriology:antibiotics}{antibiotics} (also implicated: inflammatory bowel disease,
obesity).
\end{solution}

\begin{solution}{exo:b3:microbiomes:5}
$\binom{20}{5} = \num{15504}$. Species with $10$: $1 - \binom{10}{5}/
\num{15504} = 1 - 252/\num{15504} = 0.984$; with $6$: $1 - 2002/\num{15504}
= 0.871$; with $3$: $1 - 6188/\num{15504} = 0.601$; with $1$: $1 -
\num{11628}/\num{15504} = 0.25$. $E[S_{5}] = 2.7$ species. Coverage: one
singleton, $1 - 1/20 = 0.95$.
\end{solution}

\begin{solution}{exo:b3:microbiomes:6}
$400\times 10^{11} = 4\times 10^{13}$ bacteria, about \qty{40}{g}. Genes:
$1000$ species $\times$ \num{4000}, less overlap --- of the order of
$3\times 10^{6}$ distinct genes, more than a hundred times the human
\num{20000}.
\end{solution}

\begin{solution}{exo:b3:microbiomes:7}
\omterm{def:b3:microbiomes:nodule}{Nitrogenase} is destroyed by oxygen, but the bacteroids must respire
to make sixteen ATP per N$_{2}$. The nodule cortex forms a diffusion
barrier that limits oxygen entry, and \omterm{def:b3:microbiomes:nodule}{leghaemoglobin}, at millimolar
concentration, binds the oxygen that enters so that the free
concentration stays near \qty{10}{nM} while the flux to the bacteroids
is high. Pink is oxyleghaemoglobin; a green nodule has degraded its
\omterm{def:b3:microbiomes:nodule}{leghaemoglobin} --- it is senescent and no longer fixing.
\end{solution}

\begin{solution}{exo:b3:microbiomes:8}
Both fail: no nodules and no \omterm{def:b3:microbiomes:mycorrhiza}{arbuscular mycorrhiza}, because the
calcium spike is the shared signal of the \omterm{def:b3:microbiomes:mycorrhiza}{common symbiosis pathway}.
The pathway therefore predates nodulation --- it was built for the
fungal \omterm{def:b3:microbiomes:symbiosis}{symbiosis} four hundred million years ago and recruited by the
legumes for the \omterm{def:b3:microbiomes:nodule}{rhizobia} later.
\end{solution}

\begin{solution}{exo:b3:microbiomes:9}
Kiers exposed some soybean nodules to an atmosphere of argon and
oxygen, in which the \omterm{def:b3:microbiomes:nodule}{rhizobia} could not fix nitrogen, while the rest of
the plant's nodules fixed normally; the plant cut the oxygen supply to
the non-fixing nodules and their \omterm{def:b3:microbiomes:nodule}{rhizobia} reproduced about half as
much. Without such discrimination, \omterm{def:b3:microbiomes:nodule}{rhizobia} that skipped the costly
fixation would reproduce faster than fixers, spread through the
population, and the plants would end by housing bacteria that gave
nothing: the \omterm{def:b3:microbiomes:symbiosis}{mutualism} would collapse.
\end{solution}

\begin{solution}{exo:b3:microbiomes:10}
$\binom{N-N_{i}}{n}/\binom{N}{n} = \prod_{j=0}^{n-1}(N - N_{i} - j)/(N -
j)$, each factor below $1$; going from $n$ to $n + 1$ multiplies by one
more factor below $1$, so the miss probability falls and the presence
probability rises with $n$; at $n = N$ the numerator $\binom{N -
N_{i}}{N}$ is zero and every species is present, giving $S$. For $N_{i},
n \ll N$ each factor is about $1 - N_{i}/N$, the product about $(1 -
N_{i}/N)^{n} \approx e^{-nN_{i}/N} \approx (1 - n/N)^{N_{i}}$ --- the
sampling-with-replacement approximation.
\end{solution}

\begin{solution}{exo:b3:microbiomes:11}
A male is a dead end for a symbiont that travels only in eggs, so any
trait that turns males into females, kills them in favour of their
sisters, or handicaps uninfected females raises the symbiont's
transmission. Cytoplasmic incompatibility: sperm of infected males are
modified so that the zygote dies unless the egg carries the same
\emph{Wolbachia}, which rescues it. Uninfected females therefore lose
the offspring they conceive with infected males, while infected females
lose none; the advantage grows with the infection frequency, so above a
threshold the infection spreads to fixation even if infected females
lay slightly fewer eggs --- the cost is outweighed by the fraction of
uninfected females' broods destroyed.
\end{solution}

\begin{solution}{exo:b3:microbiomes:12}
The endosymbiont passes through a bottleneck of a few cells at each
host generation and never recombines with outsiders: drift fixes
slightly deleterious mutations, deletions included, and there is no
source of replacement genes (Muller's ratchet). Genes whose products
the host supplies, or that serve a free life --- envelope, motility,
regulation, most repair --- are not missed and go first; the
translation machinery and the genes that make what the host needs
(essential amino acids in \emph{Buchnera}) are kept longest. Free-living
relatives have huge populations, recombination and gene influx, and
selection maintains their genomes. Reversal needs new genes from
outside --- impossible in isolation --- so hosts instead replace an
eroded symbiont with a fresh one, as several insect lineages have done.
\end{solution}

\begin{solution}{pb:b3:microbiomes:1}
\textbf{1.} $400\times 10^{11} = 4\times 10^{13}$ bacteria, \qty{40}{g}.
\textbf{2.} A \num{4000}-gene genome is about \qty{4}{Mb}, \qty{4.4}{fg}:
$4\times 10^{13}\times 4.4\times 10^{-15} = \qty{0.18}{g}$ of bacterial
DNA, against $3\times 10^{13}\times 6.4\times 10^{-12} = \qty{0.19}{g}$
of human DNA --- about equal.
\textbf{3.} $1000\times 0.7\times 4000 + 0.3\times 4000 \approx 2.8\times
10^{6}$ distinct genes, some $140$ times the human genome's.
\textbf{4.} $30\times 10 = \qty{300}{mmol}$; $\times 0.9 = \qty{270}{kJ}$;
\qty{3}{\%} of the diet.
\textbf{5.} No --- \qty{3}{\%} does not explain \qty{30}{\%}. Germ-free
mice also absorb less efficiently, have altered gut hormones, fat
storage and thermogenesis, and different appetite: the \omterm{def:b3:microbiomes:symbiosis}{microbiome} acts
on the host's physiology, not only as a fuel supplier.
\textbf{6.} Division balanced by voiding: a steady population, each cell
replaced daily. After a \qty{99.9}{\%} kill, $4\times 10^{10}$ remain
and ten doublings restore the numbers within about ten days --- but
the survivors and the fastest growers dominate, the composition is
changed (a \omterm{def:b3:microbiomes:gut}{dysbiosis}), and in the gap an invader such as \emph{C.
difficile} can establish.
\textbf{7.} Minor species share $0.4$ over $177$: $p = 0.00226$ each. $H
= -(0.3\ln 0.3 + 0.2\ln 0.2 + 0.1\ln 0.1 + 0.4\ln 0.00226) = 0.361 +
0.322 + 0.230 + 2.437 = 3.35$; $e^{H} \approx 28$.
\textbf{8.} $D = 0.09 + 0.04 + 0.01 + 177\times (0.00226)^{2} = 0.141$;
$1/D = 7.1$. The \omterm{def:b3:microbiomes:diversity}{Simpson index} squares the abundances, so the three
dominant species decide it and the rare ones hardly count; Shannon
gives rare species more weight.
\textbf{9.} $1 - 60/2000 = 0.97$: about $30$ reads per thousand belong
to species not yet seen.
\textbf{10.} Richness rises with sample size while rare species keep
appearing; \num{10000} reads reach species that \num{2000} miss. Rarefy
the large sample to \num{2000} reads with the theorem (or by
subsampling) --- it should give about $180$ if it is the same
community --- and compare at equal depth.
\textbf{11.} $N_{i} = 1$: $1 - \binom{1999}{1000}/\binom{2000}{1000} =
1 - (2000 - 1000)/2000 = 0.5$. $N_{i} = 5$: about $1 - 0.5^{5} = 0.97$.
\textbf{12.} $e^{2.0} = 7.4$ effective species; $D \ge 0.5^{2} = 0.25$,
about $0.27$: a community with one dominant species and a few dozen
minor ones --- the collapsed, low-diversity gut in which
\emph{C.~difficile} takes hold.
\textbf{13.} $150\times 8 = \qty{1200}{kg}$ of sugar per hectare,
\qty{10}{\%} of the \qty{12}{t} photosynthesised.
\textbf{14.} $150\,000/28 = 5360$ mol of N$_{2}$; $16\times 5360 =
\num{85700}$ mol of ATP.
\textbf{15.} $\num{85700}/30 = 2860$ mol of glucose, \qty{514}{kg} ---
about \qty{43}{\%} of the sugar paid. The rest builds and maintains the
nodules and bacteroids, supplies the eight electrons per N$_{2}$ as
reductant, and assimilates the ammonia.
\textbf{16.} $1200\times 16 = \qty{19200}{MJ}$ for \qty{150}{kg} of N:
\qty{128}{MJ/kg}, nearly four times the factory's \qty{35}{MJ/kg} --- but
paid in sunlight, in the field, with no fuel.
\textbf{17.} Fixers $0.9\times 1 = 0.9$; cheaters $0.1\times 0.5 =
0.05$; cheaters are $0.05/0.95 = 5.3\,\%$ --- halved in one season.
\textbf{18.} Cheaters $0.1\times 1.2^{10} = 0.62$ against fixers $0.9$:
\qty{41}{\%} after ten seasons, rising toward domination. Sanctions turn
selection against the cheater; without them the \omterm{def:b3:microbiomes:symbiosis}{mutualism} erodes.
\textbf{19.} $10^{6}\times \qty{20}{pg} = \qty{20}{\micro g}$ of carbon
per cm$^{2}$ per day; \qty{18}{\micro g} to the coral.
\textbf{20.} Respiration \qty{15}{\micro g}: surplus \qty{3}{\micro g}
per cm$^{2}$ per day for growth, the organic matrix of calcification,
mucus and gametes.
\textbf{21.} With \qty{5}{\%} of the algae, \qty{0.9}{\micro g} a day
against \qty{15}{\micro g} respired: a deficit of \qty{14}{\micro g} a
day; \qty{300}{\micro g} of reserves last about three weeks.
\textbf{22.} Four weeks at $+2$ is worse: the damage to the algae's
photosystems rises steeply, not linearly, with temperature. The metric
nevertheless works as a warning because bleaching integrates cumulative
stress and the two profiles differ by less than the uncertainty in the
threshold; it is an empirical index, calibrated on observed bleaching.
\textbf{23.} A stressed coral can shuffle toward more heat-tolerant algal
species already present at low abundance, or acquire them from the
water after bleaching; the tolerant algae give less carbon, so the
coral grows more slowly, and the tolerance gained is one or two
degrees --- less than the warming projected --- while the pace of
warming outruns the corals' own adaptation.
\textbf{24.} $10^{10}$ cm$^{2}$ $\times$ \qty{20}{\micro g} $=
\qty{200}{kg}$ of carbon a day, about \qty{73}{t} a year.
\textbf{25.} About \qty{3}{\%} of the diet from the \omterm{def:b3:microbiomes:symbiosis}{microbiome}'s
fermentation; \qty{97}{\%} coverage of the sequencing sample;
\qty{10}{\%} of photosynthesis paid for the bean crop's nitrogen.
\end{solution}
